% =============================================================================
%  §5 与参考文献
% =============================================================================
\section{等价}\label{sec:equiv}

我们考虑函子
\[
F_{\IW}:\;D^b\Coh^{G^\vee}(\tN)\longrightarrow D^b_{\IW}(\Fl_G),
\]
它定义为如下复合：
\[
D^b\Coh^{G^\vee}(\tN)\xrightarrow{\ F\ }D^bP_I
\longrightarrow D^b_c(I\backslash\Fl_G)\xrightarrow{\ \AvIW\ }D^b_{\IW}(\Fl_G).
\]
或者，由定理 4.4.1 与推论 4.3.2，它也可以定义为如下复合：
\[
D^b\Coh^{G^\vee}(\tN)\xrightarrow{\ F\ }D^bP_I
\xrightarrow{\ \AvIW\ }D^b\PIW\cong D^b_{\IW}(\Fl_G).
\]
可以验证 $F_{\IW}$ 对对象的映射如下：
\begin{equation}
\begin{array}{ccc}
D^b\Coh(\tN/G^\vee) & \xrightarrow{\ \bar F\ } D^b(I\backslash\Fl_G) & \xrightarrow{\ \AvIW\ } D^b_{\IW}(\Fl_G)\\[0.9em]
V\otimes\cO_{\tN} & Z\bigl(\Sat(L_\lambda)\bigr) & \mathcal{T}^{\IW}(V)\\[1.5em]
\cO_{\tN}(\lambda) & \mathcal{W}_\lambda
  & \AvIW(\mathcal{W}_\lambda)\ \cong_{K_0}\ \Delta^{\IW}_\lambda
\end{array}
\end{equation}

\begin{yizhu}[原文疑误：图 (7) 中的 $L_\lambda$]
图 (7) 第一行左端是 $V\otimes\cO_{\tN}$、右端是 $\mathcal{T}^{\IW}(V)$，中列却写成 $Z(\Sat(L_\lambda))$，
下标 $\lambda$ 在此行中并无出处，应为 $Z(\Sat(V))$。
（下一段正文「$F_{\IW}(L_\lambda\otimes\cO_{\tN})$ 是倾斜的」中的 $L_\lambda$ 另有所指：
它是最高权为 $\lambda$ 的单 $G^\vee$-模，与图中这一行的 $V$ 不是同一个记号。）
译文照录原文，此处标出。
\end{yizhu}

注意，$F_{\IW}(L_\lambda\otimes\cO_{\tN})$ 是倾斜层这一事实由定理 4.5.1 推出。
我们把这个倾斜对象记作 $\mathcal{T}^{\IW}(V)$。\footnote{自然会问：对 $\lambda\in X^\vee_+$，是否有 $\mathcal{T}^{\IW}(L_\lambda)\cong\mathcal{T}^{\IW}_\lambda$？事实上确实如此，且可作为定理 5.0.1 的推论得到。}
此外，$\AvIW(\mathcal{W}_\lambda)$ 并不同构于标准对象 $\Delta^{\IW}_\lambda$，
但它们在 $D^b_{\IW}(\Fl_G)$ 的 Grothendieck 群中代表同一个类。

我们终于来到主定理——Arkhipov–Bezrukavnikov 等价。

\begin{theorem}[Arkhipov–Bezrukavnikov, \cite{ArBe}]\label{thm:main}
函子 $F_{\IW}:D^b\Coh^{G^\vee}(\tN)\to D^b_{\IW}(\Fl_G)$ 是范畴的等价。
\end{theorem}

\begin{proof}[定理 5.0.1 的证明梗概]
$F_{\IW}$ 本质满这一事实，由 (7) 中给出的描述以及一个例行公事的同调代数论证推出。

为证明对 $\cF,\cG\in D^b\Coh^{G^\vee}(\tN)$，态射
\begin{equation}
\mathrm{Hom}_{D^b\Coh^{G^\vee}(\tN)}(\cF,\cG)
\longrightarrow
\mathrm{Hom}_{D^b_{\IW}(\Fl_G)}\bigl(F_{\IW}(\cF),F_{\IW}(\cG)\bigr)
\end{equation}
是同构，由五引理论证，只需考虑 $\cF=\cO_{\tN}$、
$\cG=V\otimes\cO_{\tN}(\lambda)[n]$ 的情形，其中 $V\in\Rep(G^\vee)$、$\lambda\in X^\vee_+$、$n\in\bbZ$。
(8) 是单射这一事实大多是直接的、并且是表示论性的。
此外，不难验证
\[
\dim\mathrm{Hom}_{D^b\Coh^{G^\vee}(\tN)}\bigl(\cO_{\tN},V\otimes\cO_{\tN}(\lambda)[n]\bigr)
=\begin{cases}
\dim V_{-\lambda}, & n=0,\\
0, & n\neq0.
\end{cases}
\]
注意存在同构
\[
F_{\IW}(\cO_{\tN})=\Delta^{\IW}_0,
\qquad
F_{\IW}\bigl(V\otimes\cO_{\tN}(\lambda)\bigr)\cong Z^{\IW}(V)\conv^{I}\mathcal{W}_\lambda .
\]
在 (8) 的右端，有
\begin{align}
(2.5.2)\qquad
&\mathrm{Hom}_{D^b_{\IW}(\Fl_G)}\bigl(\Delta^{\IW}_0,Z^{\IW}(V)\conv^{I}\mathcal{W}_\lambda[n]\bigr)
\notag\\
&\qquad\cong
\mathrm{Hom}_{D^b_{\IW}(\Fl_G)}\bigl(\Delta^{\IW}_0\conv^{I}\mathcal{W}_{-\lambda},Z^{\IW}(V)[n]\bigr)
\notag\\
(6)\qquad
&\qquad\cong
\mathrm{Hom}_{D^b_{\IW}(\Fl_G)}\bigl(\Delta^{\IW}_{-\lambda},Z^{\IW}(V)[n]\bigr).
\notag
\end{align}
于是由定理 4.5.1 可得
\[
\dim\mathrm{Hom}_{D^b_{\IW}(\Fl_G)}\bigl(\Delta^{\IW}_0,Z^{\IW}(V)\conv^{I}\mathcal{W}_\lambda[n]\bigr)
=\begin{cases}
\dim V_{-\lambda}, & n=0,\\
0, & n\neq0.
\end{cases}
\]
因此 $F_{\IW}$ 是全忠实的。
\end{proof}

\begin{yizhu}[原文含糊：证明里的两处「引证标记」]
上式左侧的 (2.5.2) 与 (6) 不是公式编号，而是指「用到命题 2.5.2」与「用到式 (6)」：
第一行用命题 2.5.2 (3) 把 $\mathcal{W}_\lambda$ 的卷积挪到左边（$\mathcal{W}_\lambda\conv\mathcal{W}_{-\lambda}\cong\mathcal{W}_0$），
第二行用式 (6) 把 $\mathrm{Av}_{\IW}(\mathcal{W}_{-\lambda})$ 换成标准层 $\Delta^{\IW}_{-\lambda}$。
原文把引证放在行首，容易误读为公式号，这里照录并说明。

另有一处需指出：原文在此处写「On the left-hand side of (8)」，但紧接着算的是
$\mathrm{Hom}_{D^b_{\IW}(\Fl_G)}\bigl(F_{\IW}(\cF),F_{\IW}(\cG)\bigr)$，
即 (8) 箭头右侧的 $\mathrm{Hom}$ 项，而非左侧的 $\mathrm{Hom}_{D^b\Coh^{G^\vee}(\tN)}$ 项。
译文据实改为「右端」，未改原文以此为准的推理。
\end{yizhu}

\begin{corollary}[Arkhipov–Bezrukavnikov, \cite{ArBe}]\label{cor:five}
存在范畴的等价
\[
\begin{array}{ccccccc}
D^b\Coh^{G^\vee}(\tN) & \xrightarrow{\ \sim\ } & D^b\Pasph
  & \xrightarrow{\ \sim\ } & D^b\PIW & \xrightarrow{\ \sim\ } & D^b_{\IW}(\Fl_G)\\[-0.2em]
 & \searrow & & \uparrow & & \nearrow & \\
 & & \multicolumn{1}{c}{D^bP_I} & & & &
\end{array}
\]
\end{corollary}

\begin{yizhu}[术语背景：五个范畴的关系]
推论 5.0.2 把本文的主要结论串成一条链：
$\tN$ 上的等变凝聚层 $\to$ 反球 Hecke 范畴的导出范畴 $\to$ Iwahori–Whittaker 范畴的导出范畴 $\to$ 平展层侧的实现。
中间夹着的 $D^bP_I$（未取反球商的仿射 Hecke 范畴）是「中转站」：
$\Pasph$ 是它的 Serre 商，$\PIW$ 是它经 $\AvIW$ 的像，
两条路最后在 $D^b_{\IW}(\Fl_G)$ 汇合（推论 4.3.2 的实现等价）。
这就是定理 1.0.1 中三个等价的具体含义。
\end{yizhu}

\begin{keytakeaways}
第 5 节的证明只有两步：先用图 (7) 得到本质满，再用五引理把全忠实化归到
$\cF=\cO_{\tN}$、$\cG=V\otimes\cO_{\tN}(\lambda)[n]$ 这一族「生成元」上，
而后者由定理 4.5.1 的重数公式直接给出。
换句话说，前面四节造机器，第五节只做一次「生成元上的比对」——这正是「等价」证明的典型形状。
\end{keytakeaways}

% =============================================================================
\section*{参考文献}
\addcontentsline{toc}{section}{参考文献}

\begin{thebibliography}{9}
\small

\bibitem[AcRi]{AcRi}
P.~Achar and S.~Riche,
\emph{Central Sheaves on Affine Flag Varieties},
可在线获取：\url{https://lmbp.uca.fr/riche/central.pdf}.

\bibitem[ArBe]{ArBe}
S.~M. Arkhipov and R.~Bezrukavnikov,
Perverse sheaves on affine flags and Langlands dual group,
\emph{Israel J. Math.} \textbf{170} (2009), 135--183.

\bibitem[BNP]{BNP}
D.~Ben-Zvi, D.~E. Nadler and A.~Preygel,
A spectral incarnation of affine character sheaves,
\emph{Compos. Math.} \textbf{153} (2017), no.~9, 1908--1944.

\bibitem[Gai1]{Gai1}
D.~Gaitsgory,
Construction of central elements in the affine Hecke algebra via nearby cycles,
\emph{Invent. Math.} \textbf{144} (2001), no.~2, 253--280.

\bibitem[Gai2]{Gai2}
D.~Gaitsgory,
The local and global versions of the Whittaker category,
\emph{Pure Appl. Math. Q.} \textbf{16} (2020), no.~3, 775--904.

\end{thebibliography}

\vfill
\noindent\rule{\linewidth}{0.4pt}
{\footnotesize
\textbf{底本}：Colton Sandvik, \emph{The Arkhipov–Bezrukavnikov Equivalence},
BunG Seminar, 2025-01-29，13 页。
本文为中文学术译注本，非官方译本，仅供研读；原文版权归原作者所有。
译注与「本节重点」为译者所加，不属于原文内容。
若关闭译注，可在导言区把 \texttt{\textbackslash showyizhutrue} 改为 \texttt{\textbackslash showyizhufalse} 重新编译。\par}
